% =====================================================================
% tesi/capitoli/23-fonti.tex
% =====================================================================
% copre: SOURCES.md
% copre: docs/24-fonti-di-community.md
% ---------------------------------------------------------------------

\chapter{Il trattamento delle fonti}
\label{cap:fonti}

La gerarchia delle fonti enunciata nella premessa non è un formalismo bibliografico: è un criterio operativo che in questo lavoro ha deciso il contenuto di sette affermazioni tecniche. Questo capitolo espone come il registro è organizzato, che cosa ciascun livello contiene, che cosa è stato effettivamente consultato e che cosa soltanto catalogato, e quali problemi di recupero si sono presentati con le fonti che non sono pagine web ordinarie.

\section{Due artefatti, due scopi}

Il registro delle fonti del progetto è un file unico condiviso da tutti i sottoprogetti. Nasce dal lavoro sul ponte fra generazioni, che è il caso di studio che ha richiesto la ricerca più profonda, ma non gli appartiene: i disassemblati, la documentazione dell'hardware, i formati di salvataggio e gli editor servono anche agli altri quattro, e tenerli in un luogo unico evita che ogni sottoprogetto riscopra le medesime cose.

Il registro elenca, dichiara a che cosa serve ciascuna voce e a quale sottoprogetto, e vale come inventario. Non dichiara però per quale ragione una fonte è stata conservata né come si lega alle altre, perché una tabella non produce un grafo. Quella parte risiede in una cartella separata, in cui ogni fonte che porta peso tecnico possiede una nota con il proprio abstract, il motivo della conservazione, il punto esatto del progetto che serve e le relazioni verso le altre. Le note sono generate da un programma a partire da una tabella unica: si modifica la tabella, non le note, e questa è la medesima disciplina del primo strato di dati generati descritta nel capitolo~\ref{cap:ponte}.

Da qui discende un vincolo: il registro è l'unica fonte di verità sulle fonti, e nessuna informazione su una fonte vive soltanto altrove. La cartella locale del materiale grezzo, esclusa dal controllo di versione, è una cache e non un archivio: contiene i collegamenti salvati durante la navigazione e i testi procurati quando il recupero automatico non funziona, cioè trascrizioni di video, pagine salvate, porzioni di discussione copiate a mano. Ciò che quel materiale documenta viene letto e trasferito nel registro in prosa, con la profondità necessaria a poterlo citare senza riaprirlo, e da quel momento il file grezzo è sacrificabile. Non si elimina per igiene, si elimina perché non porta più informazione che non sia già scritta. Il materiale grezzo resta comunque fuori dal repository, per volume e perché è contenuto di terzi: ciò che entra nel controllo di versione è la sintesi con l'attribuzione, non la copia.

Una proprietà di ogni voce merita di essere enunciata come regola di redazione, perché determina l'utilità del registro: ogni voce dichiara che cosa la fonte documenta in modo autorevole, non soltanto che essa esiste. Il valore di una fonte risiede in ciò su cui la si può citare, e un elenco di titoli non porta quella informazione.

\section{Che cosa contiene ciascun livello}

Il primo livello comprende i disassemblati, le decompilazioni e la documentazione dell'hardware, e la sua proprietà è che un dato letto in questa sede non necessita di conferma mentre un dato che lo contraddice è sbagliato. Vi appartengono i disassemblati delle tre generazioni \cite{pokered,pokecrystal,pokeemerald,pokeruby,pokefirered}, la documentazione dell'hardware del Game Boy \cite{pandocs} e del Game Boy Advance \cite{gbatek}, e la descrizione architetturale delle due macchine \cite{copetti}.

Fra le voci di questo livello vale segnalare una categoria di risultati che un registro tende a non conservare, cioè quelli negativi. Il confronto fra il disassemblato di Giallo e quello di Rosso e Blu ha stabilito che la macro della struttura di box e le costanti di lunghezza sono identiche, dunque il lettore della generazione 1 copre anche quel titolo senza modifiche \cite{pokeyellow}; il confronto fra Oro e Argento e Cristallo ha stabilito la medesima cosa per la generazione 2, restando diversi soltanto gli offset del salvataggio \cite{pokegold}. Sono risultati negativi e utili, perché ciascuno ha eliminato una variante di codice che si sarebbe altrimenti scritta per precauzione.

Il secondo livello comprende le wiki e i riferimenti di dominio, e per esso vale una precisazione sul modo di elencare: si registrano le pagine singole e non soltanto i domini, perché una wiki di grandi dimensioni è inutilizzabile come citazione mentre una pagina precisa è una fonte. Sono state lette sedici pagine enciclopediche \cite{bulbapedia}, ed è il livello che ha prodotto le quattro affermazioni errate documentate nella premessa e nei capitoli~\ref{cap:testo} e~\ref{cap:cifratura}. Vi appartiene anche il catalogo dei difetti di gioco \cite{glitchcity}, che respinge il recupero automatico e va letto dal proprio mirror statico.

Il terzo livello comprende le implementazioni di riferimento, cioè codice che funziona sul campo. Va letto come prova di fattibilità e come repertorio di soluzioni, mai come specifica, e la formulazione che vale come regola è che dove una di queste implementazioni compie una scelta, la scelta è sua e non del formato. Vi appartengono lo strumento di trasferimento \cite{ptgb}, la libreria di conversione \cite{pccs}, l'implementazione che opera nel verso opposto \cite{gen3togenx} e quella degli scambi in emulazione \cite{pokemongb-online}.

Su questo livello va introdotta una distinzione di categoria che il registro ha dovuto adottare per non confondere due situazioni diverse. Diverse voci sono strumenti eseguibili e non documenti: editor, librerie, software di lettura ed emulatori. Non sono stati letti e non avrebbe molto senso leggerli: si eseguono. Il momento in cui servono è quello in cui esiste un dato reale su cui puntarli. Sono dunque catalogati come strumenti da eseguire e non come fonti non lette, e la differenza non è terminologica: la prima categoria si impiega quando serve, la seconda costituisce un debito. Restano invece debito autentico i repository di codice del terzo livello che documentano una tecnica e che non sono stati aperti.

Il quarto livello comprende articoli, blog, diari di sviluppo e materiale video, ed è eccellente per comprendere il ragionamento e il contesto e non citabile per un offset. Vi appartengono anche una raccolta di risorse di reverse engineering sulla piattaforma \cite{retroreversing} e l'unico articolo scritto che descriva l'hardware del primo ponte \cite{hackaday-ponte}. La voce più densa dell'intero livello è la serie di articoli sullo sviluppo dello strumento di riferimento \cite{devlog-ptgb}, letta integralmente, dalla quale provengono la specifica dell'innesco esposta nel capitolo~\ref{cap:esecuzione} e la scoperta architetturale sul Dono Segreto.

Il quinto livello comprende forum e comunità, e risponde a domande che non possiedono una risposta scritta altrove. Una risposta in una discussione non è una fonte finché non è verificata. Fra le voci lette, tre discussioni di un forum specializzato hanno prodotto contributi che nessuna altra fonte forniva: il collaudo di un dispositivo generatore di clock per il protocollo, dal quale proviene la scoperta del precedente discusso nel capitolo~\ref{cap:opzioni}; la distinzione fra corruzione e assenza di dati su una cartuccia contraffatta; e il caso, discusso nel capitolo~\ref{cap:smeraldo}, di un editor che identifica erroneamente il gioco e colloca gli oggetti negli slot sbagliati, dal quale proviene il rilevamento automatico implementato nello strumento di diagnosi \cite{projectpokemon}.

\section{Le fonti che non sono pagine, e come si recuperano}

Alcune informazioni di cui questo lavoro ha necessità esistono documentate in un solo luogo, e quel luogo è una conversazione. Non un articolo, non una wiki, non un repository con la propria documentazione: un canale in cui qualcuno, in un pomeriggio di due anni prima, ha scritto tre messaggi che spiegano per quale ragione un approccio funzioni e un altro no. Il caso concreto che ha reso necessario affrontare il problema è la testimonianza sulle schede capaci di modalità monitor discussa nel capitolo~\ref{cap:ldn}: l'elenco ufficiale del progetto di riferimento indica tre modelli, il canale di supporto di una comunità ne indica un quarto che costa una frazione, e quella informazione non è scritta in alcun altro luogo.

Ne nasce un problema di metodo, perché una conversazione non è raggiungibile né dal crawler del modello né da una richiesta ordinaria: le pagine di quel servizio sono un'applicazione a pagina singola dietro autenticazione e il contenuto non è indicizzato. L'unica via è un'esportazione prodotta da chi è membro del server, cioè dall'utente.

\subsection{Il costo di quella via, dichiarato prima del resto}

Lo strumento maturo per l'esportazione richiede, per un canale di un server di cui si è membri e senza essere un programma registrato, il token dell'account personale, perché è l'unica credenziale che identifichi un account personale verso l'interfaccia del servizio.

Le condizioni d'uso del servizio non consentono l'accesso automatizzato con un account personale, e la conseguenza prevista è la sospensione dell'account. Va detto senza ambiguità, perché è una decisione di rischio e non un dettaglio di configurazione: la sospensione colpirebbe l'account e con esso l'appartenenza a tutti i server, non il file esportato. Il progetto non decide al posto dell'utente; ciò che il progetto impone è che il token non entri mai in un file tracciato né in una conversazione con l'agente, e che l'impiego resti puntuale su un canale e un intervallo di date anziché divenire continuativo.

Esiste una via che quel rischio non comporta, ed è la copia manuale della porzione di conversazione rilevante in un file di testo. È peggiore in efficienza e migliore in ogni altro aspetto, e per il volume di cui questo lavoro necessita, cioè alcune decine di messaggi su tre o quattro discussioni, è probabilmente la scelta corretta.

Il progetto possiede uno strumento che lavora su un file già prodotto e non dialoga con alcun servizio: converte l'esportazione in un formato leggibile conservando autore, momento, testo, allegati e indicazione delle risposte, e scarta il resto, che in un'esportazione costituisce la maggior parte del volume. I tre filtri disponibili sono il modo per fare entrare in contesto una conversazione senza farvi entrare il rumore: il filtro per parola chiave, ripetibile, il filtro per intervallo di date, e il filtro per lunghezza minima, che elimina le risposte di una parola, le quali in un canale di supporto costituiscono circa la metà dei messaggi e non portano informazione tecnica.

Il riconoscimento del formato di esportazione guarda la forma del documento e non l'estensione, perché l'estensione è la medesima per servizi diversi. Per uno dei formati esiste una complicazione che vale conoscere: il testo di un messaggio con formattazione non è una stringa ma una lista di frammenti, dove le parti evidenziate e i collegamenti sono oggetti, e il lettore li ricompone in ordine, altrimenti metà dei messaggi risulterebbe vuota.

Vale registrare che questo strumento non conosce nulla del dominio di questo progetto: legge un formato di esportazione e produce testo filtrato. È dunque candidato a uscire da questo progetto ed entrare nel modello di partenza come strumento disponibile a qualunque lavoro nuovo. Il passaggio va compiuto trasferendo anche la convenzione di consegna e l'avvertenza sul token, che privata del proprio contesto diverrebbe un invito implicito a violare le condizioni d'uso di un servizio.

\subsection{Il caso del forum non raggiungibile}

Un secondo servizio pone un problema analogo con esito diverso, e la sua documentazione ha valore perché la sua indisponibilità è facile da attribuire alla causa sbagliata. Non è un limite del progetto: è la somma di due fatti indipendenti, cioè che il dominio blocca il crawler del modello e che l'accesso anonimo programmatico è chiuso anche alle altre vie. Sono state tentate e documentate sei strade, tutte fallite.

Le due vie che funzionerebbero sono l'automazione del browser reale dell'utente e l'interfaccia ufficiale con credenziali applicative, per la quale esiste uno strumento pronto. Sulle condizioni di quell'interfaccia il progetto ha verificato direttamente la documentazione ufficiale, e vale registrare che essa correggeva una supposizione: l'uso gratuito esiste e non richiede alcuna approvazione preventiva per il caso non commerciale, e il modulo di contatto che la pagina cita serve alle richieste commerciali, aziendali, accademiche o di superamento dei limiti. Il limite dell'accesso gratuito è di cento richieste al minuto per identificativo, mediato su una finestra di dieci minuti così da tollerare le raffiche. È obbligatorio autenticarsi, perché il traffico non autenticato viene bloccato, ed è obbligatorio un identificativo di applicazione descrittivo, perché quelli generici delle librerie sono deliberatamente limitati. Esiste infine un obbligo che riguarda chi archivia e non chi legge: il contenuto cancellato va rimosso anche dalle copie in proprio possesso, il che costituisce una ragione ulteriore per conservare la sintesi con l'attribuzione anziché l'archivio dei messaggi.

Sul campo, tuttavia, la registrazione dell'applicazione è stata rifiutata anche dopo la verifica dell'indirizzo di posta e l'accertamento che nessun modulo preventivo fosse richiesto, con un rifiuto lato servizio che rimanda a una politica generale. La conclusione onesta è che quella via non è disponibile su quell'account e che non vale insistere.

\subsection{La forma di consegna che vale più delle altre}

Dall'esperienza di entrambi i casi è emersa una prescrizione operativa che vale registrare perché non è ovvia e perché ha prodotto il materiale migliore dell'intero lavoro di ricerca. La forma di consegna più efficace non è l'esportazione del canale, che è voluminosa e comporta il rischio descritto, e non è la copia di una conversazione scelta a occhio: è la ricerca interna al canale con un termine concordato, di cui si consegna l'elenco dei risultati.

La prova sul campo è netta. Cinque filtri su un server e tre su un altro hanno prodotto, in poche decine di schermate, il contenuto che ha corretto tre affermazioni errate del progetto e ne ha aggiunte due nuove, con un rapporto fra segnale e volume che nessuna esportazione integrale avrebbe raggiunto. La ragione è che il filtro incorpora la domanda, e chi cerca sa che cosa sta cercando.

Ne seguono tre prescrizioni per chi richieda il materiale. Si concordano i termini di ricerca prima, non il canale. Si chiede il conteggio dei risultati insieme alle schermate, perché sapere che un filtro ha prodotto zero risultati è esso stesso un dato: nel caso concreto, il filtro sul numero di un chip che ha prodotto zero risultati ha dimostrato l'assenza di qualunque testimonianza su quel componente. E si accetta che su un filtro molto generico chi cerca si fermi a una parte dei risultati, purché dichiari dove si è fermato, cosicché la copertura parziale resti dichiarata anziché apparire completa.

Il bilancio comparativo fra i due servizi va registrato perché è controintuitivo e ha una spiegazione strutturale. La resa del forum è stata molto inferiore a quella dei canali di conversazione, e la ragione è che il primo è un luogo di supporto generalista in cui la ricerca a testo pieno pesca titoli di richieste di aiuto, mentre i secondi sono luoghi in cui si discute il codice. Su tre filtri su cinque la resa è stata nulla.

\subsection{Il materiale video, e la via che funziona}

Alcuni canali video sono l'unica documentazione esistente di certe tecniche, perché i loro autori pubblicano in video ciò che non hanno mai scritto. Il limite del formato è che un video non è citabile per un offset e non è confrontabile riga per riga.

La via di recupero è stata individuata e va registrata perché non è quella immediata. La pagina del video si scarica e contiene il riferimento alla traccia dei sottotitoli automatici, ma l'endpoint che la serve restituisce zero byte a qualunque richiesta che non provenga dal lettore reale. La via praticabile è uno strumento dedicato che gestisce il token di origine e scarica i sottotitoli senza toccare l'audio; il testo si ripulisce con un programma proprio, perché i sottotitoli a scorrimento ripetono ogni riga due o tre volte. Quando i sottotitoli automatici non esistono si passa al riconoscimento vocale locale, che costa considerevolmente più tempo e resta la seconda scelta.

Sei video sono stati trascritti e letti integralmente \cite{video-goppier,video-trascritti}, e il bilancio vale registrarlo perché è anch'esso controintuitivo: due di essi sono, su punti specifici, la migliore fonte che il progetto possieda, e uno degli altri non ha prodotto nulla. Il più denso dell'intero quarto livello documenta sette circostanze che nessuna altra fonte riporta, fra cui il vincolo di sincronizzazione esposto nel capitolo~\ref{cap:cavo} e l'affermazione sulla comunicazione diretta discussa nel capitolo~\ref{cap:opzioni}. Resta arretrato il video precedente del medesimo autore, privo di sottotitoli di qualunque tipo, che è registrato come il più importante degli arretrati perché il successivo vi si riferisce più volte.

\section{Che cosa è stato usato davvero}

Un registro elenca ciò che esiste, non ciò che è stato consultato, e la differenza va dichiarata perché altrimenti l'ampiezza dell'elenco si scambia per profondità della verifica. Questa sezione è la dichiarazione.

Del primo livello sono stati clonati e cercati nel codice sei repository. Va corretta un'affermazione precedente che diceva il contrario su due voci: la documentazione dell'hardware del Game Boy e quella del Game Boy Advance sono state lette nelle pagine specifiche citate, cioè quella sul trasferimento seriale e quella sul multiboot, e non integralmente; le affermazioni tecniche di questo lavoro che vi poggiano sono quelle di quelle pagine e non altre. Gli altri repository del primo livello sono catalogati e non aperti.

Del secondo livello sono state lette le sedici pagine elencate, ed è il livello che ha prodotto le quattro affermazioni errate. Del terzo sono stati letti i sorgenti di sei progetti, mentre gli editor e le librerie sono catalogati e non eseguiti, per la ragione semplice che non esiste ancora un salvataggio reale su cui puntarli; il confronto dell'interpretazione dei campi con l'editor di riferimento è il prossimo controllo in ordine di convenienza. Del quarto livello sono stati letti gli articoli che il registro cita per contenuto specifico. Del quinto sono state lette le discussioni indicate come lette, mentre le rimanenti valgono come indicazione di dove cercare.

Nessuno strumento hardware è stato impiegato, perché nessuno è disponibile: il lettore di cartucce non è ancora giunto, e nessun emulatore è stato eseguito, nemmeno quello di cui è documentata la possibilità di collaudo attraverso una connessione TCP.


\section{Il debito di lettura, dichiarato voce per voce}

Resta un insieme di fonti catalogate e non aperte, e la loro elencazione esplicita è l'unico modo di impedire che il debito si dimentichi. Non appartengono alla categoria degli strumenti da eseguire discussa sopra: sono repository di codice o discussioni che documentano una tecnica, dunque leggerli produrrebbe conoscenza.

Sul lato del microcontrollore che dialoga con un Game Boy esistono tre precedenti non letti \cite{arduino-boy,arduino-poke-gen2,gba-link-connection}, e il primo e il secondo sono direttamente pertinenti all'opzione discussa nel capitolo~\ref{cap:opzioni}, mentre il terzo è una libreria per il collegamento dal lato del Game Boy Advance. La loro lettura è la prima cosa da compiere qualora quella opzione venga scelta, perché il capitolo~\ref{cap:opzioni} mostra che il costo stimato di una strada cambia quando si scopre che qualcuno l'ha già percorsa.

Sull'esecuzione di codice esiste un precedente non letto che opera su una piattaforma diversa da quelle trattate qui \cite{stadium-ace}, e la sua rilevanza è metodologica: documenta la medesima classe di tecnica su un sistema con vincoli differenti. Sulle distribuzioni di esemplari della terza generazione esiste un repository del medesimo autore del primo ponte \cite{gen3distributions}, non aperto, che potrebbe contenere strutture valide impiegabili come casi di prova.

Sul lato del protocollo wireless esistono due strumenti non letti che affrontano il medesimo problema del capitolo~\ref{cap:ldn} per vie diverse \cite{ldn-mitm,switch-lan-play}, cioè l'intercettazione sulla console e l'inoltro del traffico locale attraverso la rete, e un progetto più ampio nella medesima area \cite{reon}. La loro lettura non è necessaria al caso di studio come è formulato, e diverrebbe rilevante qualora l'obiettivo si estendesse oltre lo scambio locale.

Esiste infine una discussione parzialmente letta sulla correzione dei salvataggi dei titoli distribuiti come applicazioni di sistema \cite{gbatemp-vc-save}, che va registrata per ciò che è: riguarda i giochi iniettati su console e non le cartucce fisiche, dunque non risponde alla domanda del caso di studio per cui era stata cercata. La sua utilità è per il caso di studio del capitolo~\ref{cap:3ds}, qualora un giorno si iniettassero copie proprie. La discussione stessa mostra i propri limiti, perché un partecipante riporta che la sequenza di byte da individuare non esiste nella versione che possiede, e perché gli offset della seconda parte della procedura non sono mai stati individuati per due dei titoli in undici mesi di richieste.

\section{Che cosa non entra nel registro}

Non entrano le fonti che documentano come ottenere materiale escluso dal perimetro dichiarato nel capitolo~\ref{cap:perimetro}, né i salvataggi di terze parti, e non entra nulla che contenga o distribuisca materiale di chiave specifico di una console. Non entrano le pagine effimere, cioè discussioni che valgono per una singola sessione. Non entrano i mirror di ROM, che non servono ad alcuno degli obiettivi di questo lavoro, dato che si opera su cartucce possedute. E non entrano le pagine di codici di alterazione, che sono la causa più probabile del problema che il caso di studio del capitolo~\ref{cap:smeraldo} deve risolvere, e non la sua soluzione.
